\section{Complete finite-source acquisition of the signed logarithm}
\label{sec:signed-acquisition}

The local theorem can be assembled into a complete one-sided Green
estimate. We give the full construction, rather than regard a uniformly
accurate approximation as an evaluated exclusion
\cite{OGAssemblyProof,OGAssemblyAssessment}.

\subsection{The literal source and its two independent tails}
For an integer \(N\ge1\), reflect the first \(N\) positive-half-line
terms evenly and set
\[
 \Phi_N^{\rm ref}(u)=\ee^{|u|}\sum_{n=1}^N
 [2(n^2\rho(u))^2-3n^2\rho(u)]\ee^{-n^2\rho(u)},\qquad
 \widehat H(z)=\frac12\int_{-V}^V\Phi_N^{\rm ref}(u)\ee^{\ii zu}\dd u.
\]
Let \(\widehat F=D\widehat H\), with \(D\) as in \eqref{eq:F}, and put
\[
 \mathcal Q=\int_\R\ee^{9|u|}\rho(u)^2\ee^{-\rho(u)}\dd u,
 \qquad \mathcal Q^{>V}=\int_{|u|>V}\ee^{9|u|}\rho(u)^2\ee^{-\rho(u)}\dd u.
\]
The arithmetic tail inside the radial cutoff and the complete source
outside it give, with \(B_0=5C_2(0)\),
\begin{equation}\label{eq:delta-H}
 \sup_{|\Im z|\le8}|H_0(z)-\widehat H(z)|
 \le\delta_H=\tfrac12[5C_2(N)\mathcal Q+B_0\mathcal Q^{>V}].
\end{equation}
This follows immediately from Lemma~\ref{lem:source} and
\(|\ee^{\ii zu}|\le\ee^{8|u|}\). The reflected finite source need not
be smooth at zero; this undifferentiated integral argument does not
require it to be. The finite integral is entire in \(z\).

These errors have terminating explicit majorants. Successive term
ratios decrease, so
\[
 C_2(N)\le
 \frac{(N+1)^4\ee^{-\pi((N+1)^2-1)}}
 {1-((N+2)/(N+1))^4\ee^{-\pi(2N+3)}}.
\]
With \(T_V=\pi\ee^{4V}\), substitution in the radial integral gives
\begin{equation}\label{eq:radial-majorant}
 \mathcal Q^{>V}=\frac1{2\pi^{9/4}}\int_{T_V}^\infty t^{13/4}\ee^{-t}\dd t
 \le\frac{\ee^{-T_V}}{2\pi^{9/4}}
       (T_V^4+4T_V^3+12T_V^2+24T_V+24).
\end{equation}
The same bound at \(T_V=\pi\) bounds \(\mathcal Q\). Both errors tend
to zero separately.

The normalization multiplier must also be paid. On
\(|\Re z-X|\le R\), \(0\le\Im z\le8\), one sufficient bound is
\begin{equation}\label{eq:multiplier}
 |D(z)|\le\mathcal B(X,R)
 =16\pi^{9/4}\exp\left\{\frac94+
                 \frac{\pi^2[81+(|X|+R)^2]}{192}\right\}.
\end{equation}
To derive it, use the reciprocal-gamma product
\cite[5.8.2]{DLMFReciprocalGammaOC}. For \(w=s/2\),
\[
 |\Gamma(w)^{-1}|
 \le|w|\exp\{\gamma\Re w+\pi^2|w|^2/12\},
\]
because \(\log|1+w/n|\le\Re(w/n)+|w|^2/(2n^2)\).
Now \(1/2\le\Re s\le9/2\), \(|s|\ge1/2\),
\(|s|^2\le[81+(|X|+R)^2]/4\), and \(\gamma<1\).
Substitution into \eqref{eq:F} proves \eqref{eq:multiplier}.
Thus the compact error in \(F\) is at most \(\mathcal B\delta_H\).
This deliberately coarse compact bound is never integrated along an
infinite remote boundary.

\subsection{The exact boundary tail}
Use \(B,\kappa,c_0\) from Section~\ref{sec:boundary-log}, and set
\[
 L_0=\log(26/c_0),\quad N_0=2+L_0/\kappa,\quad
 C_{\log}=\log26+2B\ee^{51\kappa}N_0^2,\quad q=\ee^{-a}.
\]
Equation~\eqref{eq:canonical-absolute-log} is bounded by
\(C_{\log}[1+\log(1+|v|)]^2\) on the unit interval \(I_v\).
For an integer \(R\ge1\), covering the two remote half-lines with
unit intervals consequently gives
\begin{equation}\label{eq:assembled-log-tail}
 \int_{|x-X|>R}K(x-X)|\log|F(x)||\dd x
 \le\mathcal T(X,R):=\frac{C_{\log}}4\ee^{-aR}
 [1+\log(|X|+R+2)]^2\frac{1+q}{(1-q)^3}.
\end{equation}
Indeed the two interval weights at offset \(k\ge R\) are each at
most \(\ee^{-ak}/8\). Bound the remaining logarithmic factor by its
value at \(R\) times \((1+k-R)^2\), and sum
\(\sum_{j\ge0}q^j(1+j)^2=(1+q)/(1-q)^3\).
This is a tail of the exact logarithm, not a tail of its approximant.

\subsection{A one-sided assembly with all constants retained}
Assume \(\sup_{|x-X|\le R}|F-\widehat F|\le\delta\),
where \(0<\delta\le1\). Write
\begin{align}
 \mathcal I&=\int_{X-R}^{X+R}K(x-X)
                    \log(|\widehat F(x)|+\delta)\dd x,\nonumber\\
 \mathcal E_{X,R}(\delta)&=\sum_{k=0}^{R-1}
 \frac{\mathcal E_{X+k+1/2}(\delta)+\mathcal E_{X-k-1/2}(\delta)}
 {16\cosh(ak)},\label{eq:compact-log-envelope}\\
 U_{R,Q,\delta}(X)&=\mathcal I+\mathcal T(X,R)-P_Q(X)+\beta_Q+G(X,1).
 \label{eq:signed-upper}
\end{align}
Here \(Q\ge2\) is an integer and \(\mathcal E_v\) is exactly
\eqref{eq:canonical-envelope}.

\begin{theorem}[Complete one-sided acquisition]\label{thm:signed-acquisition}
Under these assumptions,
\begin{equation}\label{eq:signed-acquisition-error}
 0\le U_{R,Q,\delta}(X)-J_0(X)
 \le\mathcal E_{X,R}(\delta)+2\mathcal T(X,R)+2\beta_Q.
\end{equation}
For each fixed \(X\) and every positive error allowance, finite
\(R,Q,N,V\) and a positive \(\delta\) can make this bound smaller
than the allowance. No conclusion about the size of \(J_0(X)\) follows
from this acquisition assertion alone.
\end{theorem}
\begin{proof}
The integrand difference from \(\log|F|\) is nonnegative almost
everywhere. On each unit interval apply Corollary~\ref{cor:canonical-log}
and multiply by the maximum kernel weight, obtaining
\eqref{eq:compact-log-envelope}. Let \(L_{\rm comp},L_{\rm tail}\)
be the exact compact and remote logarithmic integrals. Subtracting
\eqref{eq:F-prime} from \eqref{eq:signed-upper} gives
\[
 U-J_0=(\mathcal I-L_{\rm comp})+
        (\mathcal T-L_{\rm tail})+(P-P_Q+\beta_Q).
\]
The three terms lie respectively in
\([0,\mathcal E]\), \([0,2\mathcal T]\) and \([0,2\beta_Q]\).
This proves the inequality, including both factors of two.
Given an allowance \(h>0\), choose \(R\) so that \(\mathcal T<h/6\),
then \(Q\) so that \(\beta_Q<h/6\). The finite sum \(\mathcal E\)
tends to zero with \(\delta\), so choose it with \(\mathcal E<h/3\).
Finally choose \(N,V\) using \eqref{eq:delta-H}--\eqref{eq:multiplier}
to realize that compact error. This order avoids using a multiplier
bound on an unchosen or unbounded domain.
\end{proof}

For fixed \(\delta>0\), the compact approximate integrand is continuous
and has a finite modulus of continuity. Validated subdivision can
therefore evaluate its integral to any specified positive tolerance.
If a directed evaluation satisfies \(U\le\overline U\le U+\eta_{\rm eval}\),
the complete upper error adds \(\eta_{\rm eval}\). An interval's width,
rather than an unqualified midpoint radius, pays this one-sided step.
Finite termination at each prescribed tuple is not a common all-height
complexity estimate. The next section gives a different integration
method that exploits certified local root information.
